% Generated draft: requires mathematical and source review.
\documentclass[11pt]{article}
\usepackage[margin=25mm]{geometry}
\usepackage{fontspec}
\setmainfont{Latin Modern Roman}
\usepackage{amsmath,amssymb,mathtools,unicode-math}
\setmathfont{Latin Modern Math}
\usepackage{enumitem,needspace,xurl,hyperref}
\setlength{\parindent}{0pt}
\setlength{\parskip}{5pt}
\setlength{\emergencystretch}{3em}
\newcommand{\N}{\mathbb{N}}
\newcommand{\Z}{\mathbb{Z}}
\newcommand{\Q}{\mathbb{Q}}
\newcommand{\R}{\mathbb{R}}
\newcommand{\C}{\mathbb{C}}
\newcommand{\F}{\mathbb{F}}
\newcommand{\A}{\mathbb{A}}
\newcommand{\PP}{\mathbb{P}}
\newcommand{\E}{\mathbb{E}}
\newcommand{\eps}{\varepsilon}
\newcommand{\dd}{\,\mathrm{d}}
\begin{document}
\section*{UnsolvedMath ID: 20001312}
\textbf{Source problem code:} catalogue entry
\subsection{Definitions and mathematical statement}
The assigned extraction merges Problem 22, attributed to A. Daurat and R. Gardner, with the beginning of Problem 23. The main question fixes linear forms
\[
\ell_1=x,\qquad\ell_2=y,\qquad\ell_3=2x+y,\qquad\ell_4=-x+2y
\]
and asks whether every planar convex body has an exact representation
\[
K=\{(x,y):g_1(\ell_1(x,y))+g_2(\ell_2(x,y))+g_3(\ell_3(x,y))+g_4(\ell_4(x,y))\ge0\}
\]
with measurable real univariate functions $g_i$. A ridge $g(\langle x,v\rangle)$ is constant along $v^\perp$; rescaling the two nonunit listed vectors only changes $g$. The source notes that this four-direction collection determines all planar convex bodies by X-rays and no three-direction collection does. Its broader measurable-set uniqueness claim needs the usual null-set and finite-X-ray/integrability conventions. Ridge normals and integration directions are perpendicular; the displayed collection is closed under that rotation, so its full four-direction set is unchanged.

The attached aid states that every bivariate polynomial of total degree at most three is a sum of univariate ridge polynomials on these four forms, and every quadratic uses only the first three. For an ellipse, put $\xi=x-h$, $\eta=y-k$, and $Q=a\xi^2+2b\xi\eta+c\eta^2$ with $a,c>0$, $ac>b^2$. The explicit identity
\[
Q=(a-2b)\xi^2+(c-b/2)\eta^2+(b/2)(2\xi+\eta)^2
\]
gives a three-ridge threshold $1-Q\ge0$. The univariate coefficients need not all have one sign; measurability imposes no such restriction.

The aid's complementary exact obstruction is a finite signed measure $\sigma=\sum_jc_j\delta_{x_j}\ne0$ whose pushforward by every $\ell_i$ is zero. If all positive-coefficient locations lie in $K$ and all negative-coefficient locations lie outside it, no literal threshold representation above is possible: the summed ridge values have signed sum zero, incompatible with nonnegative values inside and strictly negative ones outside. A compatible $16$-term switching family is
\[
\sigma=\mathop{*}_{i=1}^4(\delta_0-\delta_{w_i}),\qquad0\ne w_i\in\ker\ell_i,
\]
with generic lengths chosen so its sixteen subset-sum locations are distinct. The criterion concerns exact pointwise sets; a finite-point obstruction alone does not exclude a representation up to planar null sets. No convex counterexample is provided by the summary, and the general convex-body representation question remains open in its assessment.

The appended incomplete Problem 23 defines, for $a\in S^{n-1}$ and $t>0$, slabs $\{x:\|x\|_\infty\le1,\ |\langle x,a\rangle|\le t\}$ and sections $\{x:\|x\|_\infty\le1,\ \langle x,a\rangle=t\}$, with volumes $V(a,t)$ and $A(a,t)$, and $f_m=(1,\ldots,1,0,\ldots,0)/\sqrt m$. It quotes $1\le A(a,0)\le\sqrt2$, with a lower equality at $f_1$ and a truncated upper equality direction. For the printed cube $[-1,1]^n$, the actual coordinate section has area $2^{n-1}$, so the numerical claim requires normalized area $A/2^{n-1}$ or a unit-volume cube. This normalization conflict and the missing final index are left explicit rather than imported into the ridge question.
\subsection{Short English statement}
Represent convex bodies by a threshold of four fixed ridge functions. The supplied aid proves polynomial and ellipse cases and gives an exact finite switching obstruction, while keeping pointwise and almost-everywhere claims distinct. The merged cube-section fragment has a separate normalization problem.
\subsection{Sources}
\begin{itemize}
\item Local UnsolvedMath registry, permanent ID $20001312$, source code catalogue entry; canonical \texttt{aim-convex-geometry-notes.json}, \texttt{notes[43]}.
\item American Institute of Mathematics, \emph{Fourier analytic methods in convex geometry}, recovered Problem 22 attributed to A. Daurat and R. Gardner, with appended beginning of Problem 23: \url{https://aimath.org/WWN/fourierconvex/fourierconvex.pdf}. Verification reports the concatenation and comma repairs in the original PDF; no PDF inspection is claimed here.
\item The appended section names Hensley, Ball and Oleskiewicz--Pe\l czy\'nski for section bounds but gives no complete citations. The main question supplies no literature reference.
\item Attached research attempt 1: partial result, recorded run status open, rigor proved, novelty confidence low; registry status partially solved, classification PARTIAL-PROGRESS. Its polynomial and switching results are machine generated and not peer reviewed; independently verify before citation.
\end{itemize}
\end{document}
